% ===========================================================================
%  AdaFoB -- Method section
%  Requires in the preamble:
%     \usepackage{xcolor}
%     \newcommand{\needsnum}[1]{\textcolor{red}{[NUMBER PENDING: #1]}}
%     \newcommand{\needsrun}[1]{\textcolor{red}{[EXPERIMENT PENDING: #1]}}
%
%  NOMENCLATURE WARNING (read before editing):
%     BPPC, BCM and SPR are FoB's modules, published at CVPR 2026. They are
%     NOT contributions of this work and must never be described as such.
%     Our contribution is the geometry-adaptive background prompting module,
%     implemented as `refine` in experiments/train.py. It is referred to
%     throughout as GAP (Geometry-Adaptive Prompting).
% ===========================================================================

\section{Method}
\label{sec:method}

\subsection{Problem Setting}
\label{subsec:setting}

We follow the standard few-shot medical image segmentation (FSMIS) protocol.
Each episode provides a support set $\mathcal{S}=\{(\mathbf{I}^{s},\mathbf{M}^{s})\}$
of $K$ annotated slices of a novel class and a query image $\mathbf{I}^{q}$, and the
task is to predict the query mask $\mathbf{M}^{q}$ without updating the segmentation
backbone. We operate in the $1$-way $1$-shot regime ($K{=}1$) unless stated
otherwise, and we use a frozen Segment Anything Model (SAM, ViT-H) as the
segmenter. The learnable component is therefore a \emph{prompt generator}: a
function that maps $(\mathbf{I}^{s},\mathbf{M}^{s},\mathbf{I}^{q})$ to a set of
point prompts that steer SAM on the query image.

\subsection{Preliminaries: Background-Centric Prompting (FoB)}
\label{subsec:fob}

Our work builds directly on FoB~\cite{bo2026fob}, which reformulates SAM-based
FSMIS as background-prompt localisation. FoB comprises three published modules,
all of which we adopt \emph{unchanged}: Background Prompt Prototypes
Construction (BPPC), Background-centric Context Modeling (BCM), and
Structure-guided Prompt Refinement (SPR).

BPPC samples the support background prompts from a \emph{differential dilation
band} around the support mask,
\begin{equation}
  \mathcal{P} \;=\; \mathcal{U}\!\left(\rho(\mathbf{M}^{s},r) - \rho(\mathbf{M}^{s},r-\epsilon),\; N_{p}\right),
  \label{eq:fob_band}
\end{equation}
where $\rho(\cdot,r)$ denotes morphological dilation with kernel size $r$,
$\mathcal{U}(\mathcal{R},n)$ draws $n$ points uniformly from region
$\mathcal{R}$, and FoB sets $r{=}15$, $\epsilon{=}2$ and $N_{p}{=}N_{f}{=}10$.
Each sampled point is converted to a Gaussian heatmap and pooled into a prompt
prototype via masked average pooling. BCM propagates context between foreground
and background prototypes, and SPR refines the predicted query prompt
coordinates using two graphs over the prototype set: an adaptive graph
$\mathbf{A}^{ada}$ and a ring-topology graph $\mathbf{A}^{ring}$.

We stress two properties of Eq.~\eqref{eq:fob_band} that are frequently
mischaracterised. First, the band is \emph{mask-conformal}: it follows an
arbitrary contour and is not a circular ring, so FoB is not restricted to
convex targets. Second, the ring assumption enters FoB in the \emph{feature
space} of SPR through $\mathbf{A}^{ring}$, not in the image-space sampling.
Accurately stating this is necessary because our contribution is defined
against the residual rigidity of the sampler, not against a strawman.

\subsection{Motivation: Which Part of FoB Is Actually Rigid}
\label{subsec:motivation}

FoB's own limitations section states that the method ``does not yet support
highly irregular and thin structures'', attributing this to the fixed number
$N_{p}$ of background prompts and to the ring-shape prior~\cite{bo2026fob}. We
decompose that statement into three concrete, testable sources of rigidity.

\textbf{(R1) A globally fixed prompt budget.} $N_{p}$ is constant across
episodes, so the same ten points must describe a large convex organ and a small
elongated one. The number of prompts required to constrain SAM should instead
scale with the length of the boundary that must be constrained.

\textbf{(R2) A scale-independent offset.} The dilation radius $r{=}15$ is an
absolute pixel distance at $256\times256$ resolution, irrespective of target
size. FoB's own ablation shows segmentation quality is sensitive to this
proximity, so a fixed $r$ is over-tight for small targets and over-loose for
large ones.

\textbf{(R3) Arc-length- and curvature-agnostic sampling.} $\mathcal{U}(\cdot)$
spreads points uniformly over the band, allocating the same prompt density to
smooth and to highly curved boundary segments. For thin or branching
structures the dilation in Eq.~\eqref{eq:fob_band} additionally merges opposite
walls, so the differential band becomes topologically degenerate and no longer
separates adjacent branches.

\subsection{AdaFoB: Geometry-Adaptive Background Prompting}
\label{subsec:adafob}

Our central claim is a single one: \emph{the background-prompt budget and its
spatial allocation should be functions of target geometry rather than fixed
hyperparameters}. We instantiate this as one module, GAP, which replaces the
sampler of Eq.~\eqref{eq:fob_band} while leaving BPPC's prototype construction,
BCM and SPR untouched.

\textbf{Geometry-derived budget.} Given the support mask $\mathbf{M}^{s}$ we
extract its outer contour and compute the boundary length $L$, the area $A$,
the compactness $4\pi A / L^{2}$, and the medial-axis length. The prompt budget
is then obtained from a single interpretable hyperparameter $\nu$, the desired
number of prompts per unit boundary length:
\begin{equation}
  N_{p} \;=\; \mathrm{clip}\!\left(\left\lfloor \nu\, L \left(1 + \lambda\,\overline{|\kappa|}\right) \right\rceil,\; N_{\min},\; N_{\max}\right),
  \label{eq:budget}
\end{equation}
where $\overline{|\kappa|}$ is the mean absolute contour curvature and
$\lambda$ weights the curvature term. Because $N_{p}$ emerges from measurable
geometry rather than a learned regression onto an unobservable target, the
budget is interpretable, cannot overfit, and requires no additional
supervision.

\textbf{Scale-adaptive offset.} We replace the constant $r$ with a
size-dependent radius $r(A) = \mathrm{clip}(\alpha\sqrt{A/\pi}, r_{\min},
r_{\max})$, so prompt proximity tracks the effective radius of the target and
addresses (R2).

\textbf{Curvature-aware placement.} Points are drawn along the offset contour
by inverse-CDF sampling of a density $d(s)\propto 1+\lambda|\hat{\kappa}(s)|$
in arc-length parameterisation, subject to a minimum geodesic spacing $\delta$.
This concentrates prompts where the boundary is irregular and addresses (R3).
Each connected component of the offset region is processed independently and
receives at least one prompt, which removes the degeneracy of
Eq.~\eqref{eq:fob_band} on branching structures.

Because Eqs.~\eqref{eq:budget} and the placement rule share one
arc-length/curvature density, budget and placement are two readings of a single
mechanism rather than two independent modules. This yields the factorial
ablation of Sec.~\ref{subsec:ablations}: $\{$fixed, adaptive$\}$ budget
$\times$ $\{$uniform, curvature-aware$\}$ placement, with the crucial control
being a fixed budget set to the \emph{mean} adaptive budget, which isolates the
effect of allocation from the effect of simply spending more prompts.

% ---------------------------------------------------------------------------
% REVIEWER-FACING HONESTY NOTE -- keep until the corresponding runs exist.
% The claims in this section are architectural. None of them is supported by
% measured numbers yet. Do not submit until Sec.~\ref{sec:implementation}
% reports a cleared pipeline-validity gate and the ablation table is populated.
% ---------------------------------------------------------------------------
\needsrun{factorial budget/placement ablation; equal-mean-budget control}
